% ============================================================
%  paper49_mertens_function_en.tex
%  The Mertens Function: Möbius, Cancellation, and the
%  Riemann Hypothesis
%
%  Author : Michael Fuhrmann
%  Project: Specialist Mathematics Project, Build 137
%  Date   : 2026-03-12
% ============================================================

\documentclass[12pt,a4paper]{amsart}
\usepackage{amsmath,amssymb,amsthm}
\usepackage{mathtools}
\usepackage[utf8]{inputenc}
\usepackage[T1]{fontenc}
\usepackage{hyperref}
\usepackage{enumitem,booktabs,array}
\usepackage{geometry}
\geometry{margin=2.5cm}

% ---- Theorem environments ----------------------------------
\newtheorem{theorem}{Theorem}[section]
\newtheorem{lemma}[theorem]{Lemma}
\newtheorem{corollary}[theorem]{Corollary}
\newtheorem{proposition}[theorem]{Proposition}
\newtheorem{conjecture}[theorem]{Conjecture}
\theoremstyle{definition}
\newtheorem{definition}[theorem]{Definition}
\newtheorem{example}[theorem]{Example}
\newtheorem{remark}[theorem]{Remark}

% ---- Macros ------------------------------------------------
\newcommand{\Z}{\mathbb{Z}}
\newcommand{\N}{\mathbb{N}}
\newcommand{\R}{\mathbb{R}}
\newcommand{\C}{\mathbb{C}}
\newcommand{\Q}{\mathbb{Q}}
\newcommand{\Mob}{\mu}              % Möbius function
\newcommand{\Mert}{M}               % Mertens function
\newcommand{\Liouv}{\lambda}        % Liouville function
\newcommand{\BigO}{\mathcal{O}}
\DeclareMathOperator{\Li}{Li}
\DeclareMathOperator{\ord}{ord}

% ---- Title block -------------------------------------------
\title{The Mertens Function:\\
Möbius Cancellation, Dirichlet Series, and\\
Its Role in the Riemann Hypothesis}

\author{Michael Fuhrmann}
\address{Specialist Mathematics Project, Build 137}
\email{specialist-maths@localhost}
\date{2026-03-12}

% ============================================================
\begin{document}

\begin{abstract}
The Mertens function $M(n) = \sum_{k=1}^{n} \mu(k)$, where $\mu$ denotes
the Möbius function, occupies a central position in analytic number theory.
Its growth rate is provably equivalent to the Riemann Hypothesis: $M(n) =
\BigO(n^{1/2+\varepsilon})$ for every $\varepsilon > 0$ if and only if the
Riemann Hypothesis holds.
Three classical theorems of Mertens (1874) govern the distribution of
reciprocal primes and the Euler product up to a given bound.
The famous Mertens Conjecture $|M(n)| < \sqrt{n}$, long believed and
widely tested, was definitively \emph{disproved} in 1985 by Odlyzko and
te~Riele using the theory of zeros of the Riemann zeta function.
We survey the Möbius function, Mertens' three theorems, the Dirichlet series
$1/\zeta(s)$, the equivalence with the Riemann Hypothesis, the Liouville
function, Selberg's elementary proof technique, and computational algorithms.
\end{abstract}

\maketitle
\tableofcontents

% ============================================================
\section{Introduction}
\label{sec:intro}

The interplay between prime numbers and the integers has, since Gauss and
Dirichlet, been one of the most productive themes in mathematics.  Among the
functions that encode this relationship, the \emph{Mertens function}
\[
  M(n) \;=\; \sum_{k=1}^{n} \mu(k)
\]
is distinguished by two properties: it is elementary to define, yet its
asymptotic behaviour is equivalent to one of the deepest unsolved problems in
mathematics.

\subsection*{Historical context}

Franz Mertens (1840--1927) was an Austrian--German mathematician whose 1874
paper \cite{Mertens1874} established three fundamental estimates---now called
Mertens' First, Second, and Third Theorems---governing the partial sums of
$1/p$, $\log p/p$, and the Euler product $\prod_{p \le x}(1-1/p)$.  These
theorems provided the first quantitative understanding of how primes thin out.

In 1897, Mertens observed numerically that $|M(n)| < \sqrt{n}$ for all tested
values and asked whether this inequality might hold for all $n \ge 1$.  This
\emph{Mertens Conjecture} remained unresolved for nearly ninety years.  In
1985, Odlyzko and te~Riele \cite{OdlyzkoTeRiele1985} proved, using a
large-scale computation of zeta zeros, that the conjecture is \emph{false}.
Paradoxically, no explicit counterexample has ever been exhibited; the
smallest counterexample is believed (heuristically) to lie near
$x \approx e^{1.59 \times 10^{40}}$.

\subsection*{The connection to the Riemann Hypothesis}

Let $\zeta(s)$ denote the Riemann zeta function.  The Dirichlet series
\[
  \frac{1}{\zeta(s)} \;=\; \sum_{n=1}^{\infty} \frac{\mu(n)}{n^s},
  \qquad \operatorname{Re}(s) > 1,
\]
provides the link between $\mu$ and $\zeta$.  The Riemann Hypothesis (RH)
asserts that every non-trivial zero $\rho$ of $\zeta(s)$ satisfies
$\operatorname{Re}(\rho) = 1/2$.  A celebrated equivalence theorem states:

\begin{center}
  \textbf{RH holds $\Longleftrightarrow$ $M(x) = \BigO(x^{1/2+\varepsilon})$
  for every $\varepsilon > 0$.}
\end{center}

\subsection*{Outline}

Section~\ref{sec:mobius} reviews the Möbius function.
Section~\ref{sec:mertens-function} defines $M(n)$ and records its Dirichlet
series.
Section~\ref{sec:three-theorems} states and proves Mertens' three theorems.
Section~\ref{sec:rh} proves the equivalence of RH and the square-root bound on
$M$.
Section~\ref{sec:mertens-conjecture} discusses the Mertens Conjecture and its
disproof.
Section~\ref{sec:liouville} treats the Liouville function.
Section~\ref{sec:selberg} outlines Selberg's elementary proof of the PNT via
$M$.
Section~\ref{sec:algorithms} covers computational algorithms.

% ============================================================
\section{The Möbius Function}
\label{sec:mobius}

\subsection{Definition}

\begin{definition}[Möbius function]
\label{def:mobius}
Let $n \in \N$.  The \emph{Möbius function} $\mu : \N \to \{-1,0,1\}$ is
defined by
\[
  \mu(n) \;=\;
  \begin{cases}
    1                  & \text{if } n = 1, \\
    (-1)^{k}           & \text{if } n = p_1 p_2 \cdots p_k
                         \text{ with distinct primes } p_i, \\
    0                  & \text{if } p^2 \mid n \text{ for some prime } p.
  \end{cases}
\]
Equivalently, $\mu(n) = 0$ when $n$ is not \emph{squarefree}, and
$\mu(n) = (-1)^{\omega(n)}$ when $n$ is squarefree, where $\omega(n)$ denotes
the number of distinct prime factors.
\end{definition}

\subsection{First values}

\begin{example}
\label{ex:mobius-values}
The first twenty values of $\mu(n)$:
\begin{center}
\begin{tabular}{r|rrrrrrrrrrrrrrrrrrrr}
$n$      & 1 & 2 & 3 & 4 & 5 & 6 & 7 & 8 & 9 & 10
         & 11 & 12 & 13 & 14 & 15 & 16 & 17 & 18 & 19 & 20 \\
\hline
$\mu(n)$ & 1 & $-1$ & $-1$ & 0 & $-1$ & 1 & $-1$ & 0 & 0 & 1
         & $-1$ & 0 & $-1$ & 1 & 1 & 0 & $-1$ & 0 & $-1$ & 0
\end{tabular}
\end{center}
Note $\mu(4)=\mu(8)=\mu(9)=\mu(12)=0$ because $4=2^2$, $8=2^3$,
$9=3^2$, $12=4\cdot 3$.
\end{example}

\subsection{Multiplicativity}

\begin{theorem}[Multiplicativity of $\mu$]
\label{thm:mobius-mult}
The Möbius function is \emph{multiplicative}: if $\gcd(m,n)=1$ then
\[
  \mu(mn) = \mu(m)\,\mu(n).
\]
\end{theorem}

\begin{proof}
If either $m$ or $n$ contains a square prime factor, then so does $mn$, and
both sides are zero.  Otherwise, $m = p_1 \cdots p_j$ and $n = q_1 \cdots q_k$
with all $p_i, q_l$ distinct (since $\gcd(m,n)=1$).  Then
$mn = p_1 \cdots p_j q_1 \cdots q_k$ is squarefree with $j+k$ prime factors,
so $\mu(mn)=(-1)^{j+k}=(-1)^j(-1)^k=\mu(m)\mu(n)$.
\end{proof}

\begin{theorem}[Orthogonality]
\label{thm:mobius-ortho}
For $n \ge 1$:
\[
  \sum_{d \mid n} \mu(d) \;=\; [n = 1] \;=\;
  \begin{cases} 1 & n=1, \\ 0 & n>1. \end{cases}
\]
\end{theorem}

\begin{proof}
The function $\varepsilon(n) = [n=1]$ is the Dirichlet identity.  Writing
$\mathbf{1}(n)=1$ for all $n$, the identity $\mu * \mathbf{1} = \varepsilon$
follows from $\zeta(s)\cdot(1/\zeta(s))=1$ in the half-plane
$\operatorname{Re}(s)>1$, or by direct computation on prime powers.
\end{proof}

\subsection{Sieve interpretation}

The formula $\sum_{d \mid n}\mu(d) = [n=1]$ is the basis of \emph{Möbius
inversion}: if $f = g * \mathbf{1}$ then $g = f * \mu$.  In particular it
underlies the inclusion-exclusion principle in combinatorial sieving, where
$\mu$ acts as the signed indicator of inclusion/exclusion rounds.

% ============================================================
\section{The Mertens Function}
\label{sec:mertens-function}

\subsection{Definition and basic properties}

\begin{definition}[Mertens function]
\label{def:mertens}
The \emph{Mertens function} is the partial-sum function
\[
  M(x) \;=\; \sum_{1 \le n \le x} \mu(n), \qquad x \ge 1.
\]
Since $\mu$ takes only integer values, $M(x) = M(\lfloor x \rfloor)$ is a
step function that changes at each integer.
\end{definition}

\subsection{First values}

\begin{example}
\label{ex:mertens-values}
Using the table in Example~\ref{ex:mobius-values}:
\begin{center}
\begin{tabular}{r|rrrrrrrrrr}
$n$    & 1 & 2 & 3 & 4 & 5 & 6 & 7 & 8 & 9 & 10 \\
\hline
$M(n)$ & 1 & 0 & $-1$ & $-1$ & $-2$ & $-1$ & $-2$ & $-2$ & $-2$ & $-1$
\end{tabular}
\end{center}

Selected larger values (computed):
\[
  M(10) = -1, \quad M(100) = 1, \quad M(1000) = 2, \quad M(10000) = -23.
\]
These small values illustrate a remarkable cancellation: roughly half the
squarefree integers have an even number of prime factors and half have an odd
number.
\end{example}

\subsection{Dirichlet series representation}

\begin{theorem}[Dirichlet series of $1/\zeta$]
\label{thm:dirichlet-series}
For $\operatorname{Re}(s) > 1$,
\[
  \frac{1}{\zeta(s)} \;=\; \sum_{n=1}^{\infty} \frac{\mu(n)}{n^s}.
\]
\end{theorem}

\begin{proof}
The Euler product $\zeta(s) = \prod_p (1-p^{-s})^{-1}$ gives
\[
  \frac{1}{\zeta(s)} = \prod_p (1 - p^{-s}).
\]
Expanding the product over all primes and collecting terms with $n=p_1^{a_1}
\cdots p_k^{a_k}$: the coefficient of $n^{-s}$ is $(-1)^k$ if all exponents
$a_i=1$ (squarefree), and $0$ otherwise (since $(1-p^{-s})$ contributes at
most $-p^{-s}$).  This is exactly $\mu(n)$.
\end{proof}

\begin{remark}
The series $\sum \mu(n)/n^s$ can be analytically continued beyond
$\operatorname{Re}(s)=1$.  Its analytic behaviour is controlled by the zeros
of $\zeta(s)$: the series converges for $\operatorname{Re}(s) > 1/2$
\emph{if and only if} the Riemann Hypothesis holds.
\end{remark}

\subsection{Explicit formula}

Under standard hypotheses on the zeros of $\zeta$, a Perron-type inversion
gives the \emph{explicit formula}
\begin{equation}
\label{eq:explicit}
  M(x) \;=\; \sum_{\rho}\frac{x^{\rho}}{\rho\,\zeta'(\rho)}
           \;-\; 2 \;+\; \sum_{n=1}^{\infty} \frac{(-1)^{n+1}}{(2n)!\,(2n)\,\zeta(2n+1)}
           \cdot \frac{1}{x^{2n}} \;+\; \cdots
\end{equation}
where the sum is over non-trivial zeros $\rho$.  The dominant terms come from
zeros with $\operatorname{Re}(\rho)$ close to $1$.  Under RH all non-trivial
zeros have $\operatorname{Re}(\rho)=1/2$, bounding every term by
$x^{1/2}/|\rho\,\zeta'(\rho)|$.

% ============================================================
\section{Mertens' Theorems}
\label{sec:three-theorems}

In his 1874 paper, Mertens proved three estimates that remain cornerstones of
analytic number theory.

\subsection{First Theorem: partial sums of \texorpdfstring{$\log p/p$}{log p/p}}

\begin{theorem}[Mertens' First Theorem, 1874]
\label{thm:mertens1}
As $x \to \infty$,
\[
  \sum_{p \le x} \frac{\log p}{p} \;=\; \log x \;+\; \BigO(1).
\]
\end{theorem}

\begin{proof}[Sketch]
From the prime number theorem in the form
$\sum_{n \le x} \Lambda(n) = x + \BigO(xe^{-c\sqrt{\log x}})$ (where
$\Lambda$ is the von Mangoldt function), partial summation yields
\[
  \sum_{p \le x} \frac{\log p}{p}
  = \sum_{p \le x} \frac{\Lambda(p)}{p}
  = \log x - \sum_{p,k \ge 2} \frac{\log p}{p^k} + \BigO(1)
  = \log x + \BigO(1),
\]
since the sum over prime powers $p^k$ with $k \ge 2$ converges absolutely.
\end{proof}

\subsection{Second Theorem: partial sums of \texorpdfstring{$1/p$}{1/p}}

\begin{theorem}[Mertens' Second Theorem, 1874]
\label{thm:mertens2}
There exists a constant $M_0$ (the \emph{Meissel--Mertens constant},
$M_0 \approx 0.2614972$) such that
\[
  \sum_{p \le x} \frac{1}{p} \;=\; \log\log x \;+\; M_0 \;+\; \BigO\!\left(\frac{1}{\log x}\right).
\]
\end{theorem}

\begin{proof}[Sketch]
Apply Abel summation (partial summation) to the sum $\sum_{p \le x} 1/p$
using $\sum_{p \le t} \log p / p = \log t + c + \BigO(1/\log t)$ from
Theorem~\ref{thm:mertens1}:
\[
  \sum_{p \le x}\frac{1}{p}
  = \frac{\log\log x + c}{\log x/\log x}
  + \int_2^x \frac{\log t + c}{t\,(\log t)^2}\,dt + \BigO(1/\log x).
\]
The integral evaluates to $\log\log x + M_0' + \BigO(1/\log x)$ for an
explicit constant, giving the result.
\end{proof}

\begin{remark}
The Meissel--Mertens constant $M_0$ is defined by
\[
  M_0 \;=\; \gamma \;+\; \sum_p \!\left(\log\!\left(1-\frac{1}{p}\right) + \frac{1}{p}\right),
\]
where $\gamma \approx 0.5772156649$ is the Euler--Mascheroni constant.
\end{remark}

\subsection{Third Theorem: the Euler product over primes}

\begin{theorem}[Mertens' Third Theorem, 1874]
\label{thm:mertens3}
As $x \to \infty$,
\[
  \prod_{p \le x} \left(1 - \frac{1}{p}\right)
  \;\sim\; \frac{e^{-\gamma}}{\log x}.
\]
More precisely,
\[
  \prod_{p \le x} \left(1 - \frac{1}{p}\right)
  \;=\; \frac{e^{-\gamma}}{\log x}
        \left(1 + \BigO\!\left(\frac{1}{\log x}\right)\right).
\]
\end{theorem}

\begin{proof}[Sketch]
Take logarithms:
\[
  \log \prod_{p \le x}\!\left(1-\frac{1}{p}\right)
  = \sum_{p \le x} \log\!\left(1-\frac{1}{p}\right)
  = -\sum_{p \le x}\frac{1}{p}
    - \sum_{p \le x}\sum_{k=2}^{\infty}\frac{1}{k\,p^k}.
\]
The tail $\sum_p \sum_{k \ge 2} 1/(kp^k)$ converges absolutely to a constant
$C$.  By Theorem~\ref{thm:mertens2},
\[
  \log \prod_{p \le x}\!\left(1-\frac{1}{p}\right)
  = -\log\log x - M_0 - C + \BigO(1/\log x).
\]
One shows $M_0 + C = \gamma$, giving
$\prod_{p \le x}(1-1/p) = e^{-\gamma}/\log x \cdot (1+\BigO(1/\log x))$.
\end{proof}

\begin{remark}
Mertens' Third Theorem is closely related to the probability that a random
integer near $x$ is \emph{not} divisible by any prime $\le x$, relevant to
sieve theory and the density of primes.
\end{remark}

% ============================================================
\section{Connection to the Riemann Hypothesis}
\label{sec:rh}

\subsection{The equivalence theorem}

\begin{theorem}[RH--Mertens equivalence]
\label{thm:rh-equiv}
The following are equivalent:
\begin{enumerate}[label=(\roman*)]
  \item The Riemann Hypothesis: every non-trivial zero $\rho$ of $\zeta(s)$
        satisfies $\operatorname{Re}(\rho) = 1/2$.
  \item For every $\varepsilon > 0$: $M(x) = \BigO(x^{1/2+\varepsilon})$
        as $x \to \infty$.
  \item The Dirichlet series $\sum_{n=1}^{\infty}\mu(n)/n^s$ converges for all
        $s$ with $\operatorname{Re}(s) > 1/2$.
\end{enumerate}
\end{theorem}

\begin{proof}[Proof sketch]
\textbf{(i)$\Rightarrow$(ii).}
Perron's formula expresses $M(x)$ as a contour integral
\[
  M(x) = \frac{1}{2\pi i}\int_{c-i\infty}^{c+i\infty}
          \frac{x^s}{\zeta(s)\,s}\,ds, \qquad c > 1.
\]
Shifting the contour to $\operatorname{Re}(s)=1/2+\varepsilon$ and collecting
residues at non-trivial zeros $\rho$ of $\zeta$, each zero contributes a term
$x^{\rho}/(\rho\,\zeta'(\rho))$.  Under RH, $\operatorname{Re}(\rho)=1/2$,
so $|x^{\rho}| = x^{1/2}$.  Summing the (conditionally convergent) series of
residues gives $M(x) = \BigO(x^{1/2+\varepsilon})$.

\textbf{(ii)$\Rightarrow$(i).}
Suppose $M(x) = \BigO(x^{1/2+\varepsilon})$ for every $\varepsilon > 0$.
Integration by parts shows that the Dirichlet series
$\sum\mu(n)/n^s$ converges absolutely for $\operatorname{Re}(s)>1/2+\varepsilon$.
Hence $1/\zeta(s)$ extends analytically to the half-plane
$\operatorname{Re}(s)>1/2$, which implies $\zeta(s) \ne 0$ there.  Since
$\zeta$ is known to satisfy the functional equation, all non-trivial zeros
must lie on $\operatorname{Re}(s)=1/2$.

\textbf{(ii)$\Leftrightarrow$(iii)} follows by standard partial summation.
\end{proof}

\subsection{Pintz's unconditional result}

\begin{theorem}[Pintz, 1987]
\label{thm:pintz}
Unconditionally (i.e., without assuming RH):
\[
  \limsup_{x \to \infty} \frac{M(x)}{\sqrt{x}} \;>\; 1.06.
\]
\end{theorem}

This follows from the 1985 result of Odlyzko and te~Riele (see
Section~\ref{sec:mertens-conjecture}) and a careful refinement of their
argument.

\subsection{Conditional bounds}

Under RH, more precise estimates are available:

\begin{proposition}[Conditional bound]
\label{prop:cond-bound}
Assume RH.  Then
\[
  M(x) \;=\; \BigO\!\left(x^{1/2}\exp\!\left(C\,\frac{\sqrt{\log x}}{\log\log x}\right)\right)
\]
for an effectively computable constant $C > 0$.
\end{proposition}

% ============================================================
\section{The Mertens Conjecture and Its Disproof}
\label{sec:mertens-conjecture}

\subsection{Statement of the conjecture}

Numerical evidence suggested a much stronger bound than what RH implies:

\begin{conjecture}[Mertens Conjecture, 1897 -- \textbf{FALSE}]
\label{conj:mertens}
For all $n \ge 1$:
\[
  |M(n)| \;<\; \sqrt{n}.
\]
\end{conjecture}

This conjecture was tested extensively.  It holds for all $n \le 10^{13}$
(Kotnik and te~Riele, 2006).

\subsection{The Odlyzko--te Riele disproof}

\begin{theorem}[Odlyzko--te~Riele, 1985]
\label{thm:odlyzko}
The Mertens Conjecture is \emph{false}: there exist arbitrarily large $x$ for
which $|M(x)| \ge \sqrt{x}$.
\end{theorem}

\begin{proof}[Method (sketch)]
Odlyzko and te~Riele \cite{OdlyzkoTeRiele1985} argued as follows.  Under the
assumption that RH holds and that non-trivial zeros $\rho_n = 1/2 + i\gamma_n$
are simple, the explicit formula \eqref{eq:explicit} gives
\[
  \frac{M(x)}{\sqrt{x}} \;=\;
  2\,\Re\sum_{\rho} \frac{x^{i\gamma}}{\rho\,\zeta'(\rho)}
  \;+\; \BigO(x^{-1/2}).
\]
The key insight is that $M(x)/\sqrt{x}$ behaves like a quasi-random
trigonometric sum.  By computing a large number of zeros
($\gamma_1, \ldots, \gamma_{2000}$ at high precision) and applying the
Gram--Schmidt process and lattice reduction, they showed that
\[
  \limsup_{x\to\infty}\frac{M(x)}{\sqrt{x}} \;\ge\; 1.06
  \quad\text{and}\quad
  \liminf_{x\to\infty}\frac{M(x)}{\sqrt{x}} \;\le\; -1.009.
\]
Both exceed $1$ in absolute value, disproving Conjecture~\ref{conj:mertens}.
\end{proof}

\begin{remark}
\label{rem:counterexample}
Despite the theoretical disproof, no \emph{explicit} counterexample $n$ with
$|M(n)| \ge \sqrt{n}$ has ever been found.  Heuristic considerations predict
the smallest counterexample to be near $e^{1.59 \times 10^{40}}$.

The disproof is \emph{non-constructive}: it shows the conjecture must fail
without producing the failing value.
\end{remark}

\subsection{Why the conjecture seemed plausible}

The Dirichlet series $\sum \mu(n)/n^s = 1/\zeta(s)$ suggests that $\mu$ is
``random-like.'' A heuristic model treating $\mu(n)$ as i.i.d.\ with $P(\mu=1)
= P(\mu=-1) = 3/\pi^2 \approx 0.304$ (density of squarefree integers with
$\omega$ even, odd) and $P(\mu=0) = 1 - 6/\pi^2 \approx 0.392$ predicts
$M(n)/\sqrt{n} \to N(0,\sigma^2)$ by CLT-type arguments, consistent with
$|M(n)| < C\sqrt{n\log\log n}$ with probability 1.  The sharper Mertens bound
$|M(n)| < \sqrt{n}$ (without $\log\log n$) is then essentially a ``$1\sigma$''
bound, which the probabilistic model shows will eventually be violated.

% ============================================================
\section{The Liouville Function}
\label{sec:liouville}

\subsection{Definition}

\begin{definition}[Liouville function]
\label{def:liouville}
Let $\Omega(n)$ denote the number of prime factors of $n$ counted with
multiplicity ($\Omega(1)=0$, $\Omega(12)=\Omega(4\cdot 3)=3$).  The
\emph{Liouville function} is
\[
  \lambda(n) \;=\; (-1)^{\Omega(n)}.
\]
The \emph{summatory Liouville function} is
\[
  L(x) \;=\; \sum_{n \le x} \lambda(n).
\]
\end{definition}

\subsection{First values}

\begin{example}
\label{ex:liouville-values}
\begin{center}
\begin{tabular}{r|rrrrrrrrrr}
$n$       & 1 & 2 & 3 & 4 & 5 & 6 & 7 & 8 & 9 & 10 \\
\hline
$\Omega(n)$ & 0 & 1 & 1 & 2 & 1 & 2 & 1 & 3 & 2 & 2 \\
$\lambda(n)$& 1 &$-1$&$-1$& 1 &$-1$& 1 &$-1$&$-1$& 1 & 1
\end{tabular}
\end{center}
Note that $\lambda$ differs from $\mu$: e.g.\ $\lambda(4)=1$ while $\mu(4)=0$.
\end{example}

\subsection{Properties and Dirichlet series}

\begin{theorem}[Complete multiplicativity]
\label{thm:liouville-mult}
The Liouville function is \emph{completely multiplicative}: $\lambda(mn) =
\lambda(m)\lambda(n)$ for \emph{all} $m,n \ge 1$ (no coprimality required).
\end{theorem}

\begin{proof}
$\Omega(mn) = \Omega(m) + \Omega(n)$ always, so $\lambda(mn) = (-1)^{\Omega(m)+\Omega(n)}
= (-1)^{\Omega(m)}(-1)^{\Omega(n)} = \lambda(m)\lambda(n)$.
\end{proof}

\begin{theorem}[Dirichlet series of $\lambda$]
\label{thm:liouville-dirichlet}
For $\operatorname{Re}(s) > 1$:
\[
  \sum_{n=1}^{\infty} \frac{\lambda(n)}{n^s}
  \;=\; \frac{\zeta(2s)}{\zeta(s)}.
\]
\end{theorem}

\begin{proof}
Using the Euler product and $\lambda(p^k) = (-1)^k$:
\[
  \sum_{n=1}^{\infty}\frac{\lambda(n)}{n^s}
  = \prod_p \sum_{k=0}^{\infty}\frac{(-1)^k}{p^{ks}}
  = \prod_p \frac{1}{1+p^{-s}}
  = \prod_p \frac{1-p^{-2s}}{1-p^{-s}} \cdot \frac{1}{1-p^{-2s}}\cdot(1-p^{-2s})
  = \frac{\zeta(2s)}{\zeta(s)}.
\]
\end{proof}

\subsection{Relation to the Möbius function}

\begin{proposition}
\label{prop:lambda-mu}
For $n \ge 1$:
\[
  \mu(n) \;=\; \sum_{d^2 \mid n} \lambda\!\left(\frac{n}{d^2}\right)
  \qquad\text{and}\qquad
  \lambda(n) \;=\; \sum_{d \mid n} |\mu(n/d)|. \quad\text{(informally)}
\]
More precisely, $\mu = \lambda * \mathbf{1}_{|\mu|=1}$ (Dirichlet convolution).
\end{proposition}

\subsection{The Pólya Conjecture}

\begin{conjecture}[Pólya, 1919 -- \textbf{FALSE}]
\label{conj:polya}
$L(n) \le 0$ for all $n \ge 2$.
\end{conjecture}

\begin{theorem}[Haselgrove, 1958]
\label{thm:haselgrove}
The Pólya Conjecture is \emph{false}.  The smallest counterexample is
$n = 906316571$.
\end{theorem}

The Pólya Conjecture is the Liouville analogue of the Mertens Conjecture; it
was disproved some 25 years earlier, with an explicit counterexample.

% ============================================================
\section{Selberg's Formula and the Prime Number Theorem}
\label{sec:selberg}

\subsection{Selberg's fundamental formula}

A landmark in elementary analytic number theory is Selberg's 1949 formula,
which provides an ``almost-PNT'' statement without complex analysis.

\begin{theorem}[Selberg's Formula, 1949]
\label{thm:selberg}
\[
  \sum_{n \le x} \Lambda(n)\log n
  \;+\; \sum_{mn \le x} \Lambda(m)\Lambda(n)
  \;=\; 2x\log x + \BigO(x),
\]
where $\Lambda$ is the von Mangoldt function.
\end{theorem}

\subsection{Connection to $M(x)$}

The Mertens function enters through the identity
\[
  \sum_{n \le x} \mu(n)\left\lfloor\frac{x}{n}\right\rfloor \;=\; 1
\]
(a consequence of Theorem~\ref{thm:mobius-ortho}).  Equivalently,
\[
  \sum_{n \le x}\mu(n)\left(\frac{x}{n} + \BigO(1)\right) = 1
  \implies x\sum_{n \le x}\frac{\mu(n)}{n} = 1 + \BigO(M(x)).
\]

\begin{theorem}[PNT via $M$]
\label{thm:pnt-mertens}
The Prime Number Theorem $\pi(x) \sim x/\log x$ is equivalent to
\[
  M(x) = o(x), \qquad x \to \infty.
\]
\end{theorem}

\begin{proof}[Sketch]
One direction: PNT implies $\sum_{n \le x}\mu(n)/n \to 0$, which by partial
summation is equivalent to $M(x)/x \to 0$.
Conversely, $M(x)=o(x)$ can be fed back into Selberg's formula to deduce the
PNT.  Selberg and Erd\H{o}s (independently, 1949) completed this elementary
argument.
\end{proof}

\subsection{Sharper equivalences}

Each sharpening of the error term in the PNT corresponds to a sharpening of
the bound on $M(x)$:
\begin{align*}
  M(x) &= \BigO(x\,e^{-c\sqrt{\log x}}) &&\Longleftrightarrow&&
    \psi(x) = x + \BigO(x\,e^{-c\sqrt{\log x}}), \\
  M(x) &= \BigO(x^{1/2+\varepsilon}) &&\Longleftrightarrow&&
    \text{Riemann Hypothesis.}
\end{align*}

% ============================================================
\section{Computational Aspects}
\label{sec:algorithms}

\subsection{Naive algorithm}

The straightforward approach computes $\mu(n)$ for each $n$ by trial division
and accumulates the sum.  Cost: $\BigO(n\sqrt{n})$ time, $\BigO(1)$ space
(streaming).

\subsection{Sieve algorithm (Eratosthenes-style)}

\begin{enumerate}[label=\arabic*.]
  \item Initialise array \texttt{mu[1..N]} with \texttt{mu[1] = 1}.
  \item Use a \emph{linear sieve} (smallest-prime-factor sieve) to compute
        $\mu(n)$ for all $n \le N$:
        \begin{itemize}
          \item For each prime $p$: set \texttt{mu[p] = -1}.
          \item For composite $n = p \cdot m$: if $p \mid m$ then
                \texttt{mu[n] = 0}; otherwise \texttt{mu[n] = -mu[m]}.
        \end{itemize}
  \item Compute prefix sums: \texttt{M[n] = M[n-1] + mu[n]}.
\end{enumerate}

Time: $\BigO(N \log\log N)$.  Space: $\BigO(N)$.  This is the algorithm
implemented in the \texttt{MertensFunction} class of the companion module
\texttt{mertens\_function.py}.

\subsection{Sub-linear algorithm}

For very large $N$, storing the full sieve is impractical.  The
Meissel--Lehmer method (Lucy\_Hedgehog variant) computes $M(N)$ in
$\BigO(N^{2/3})$ time and $\BigO(N^{1/3})$ space using the identity
\[
  M(N) \;=\; 1 \;-\; \sum_p M\!\left(\left\lfloor\frac{N}{p}\right\rfloor\right)
            \;-\; \text{(correction terms)}.
\]
This approach is based on the Dirichlet hyperbola method applied to $\mu$.

\begin{theorem}[Sub-linear Mertens computation]
\label{thm:sublinear}
$M(N)$ can be computed in $\BigO(N^{2/3})$ arithmetic operations and
$\BigO(N^{1/3})$ memory.
\end{theorem}

\subsection{Known values and numerical records}

\begin{center}
\begin{tabular}{rl}
\toprule
$N$ & $M(N)$ \\
\midrule
$10^4$  & $-23$ \\
$10^6$  & $212$ \\
$10^8$  & $754$ \\
$10^{10}$ & $-1215$ \\
$10^{12}$ & $-2$\\
$10^{13}$ & $-3\,994\,941$ \\
\bottomrule
\end{tabular}
\end{center}

The ratio $M(N)/\sqrt{N}$ remains well below $1$ for all computed values;
the largest known value is $|M(N)/\sqrt{N}| \approx 0.571$ at
$N \approx 7.76 \times 10^9$.

\subsection{Computational complexity and RH}

If RH holds, efficient algorithms for $M(N)$ can be designed exploiting
the $\BigO(N^{1/2+\varepsilon})$ bound.  Conversely, a hypothetical algorithm
that computes $M(N)$ faster than $\BigO(N^{1/2})$ would provide evidence
(though not proof) for RH.

% ============================================================
\section*{Conclusion}

The Mertens function sits at the crossroads of three fundamental themes:
\begin{enumerate}[label=(\roman*)]
  \item \textbf{Multiplicative number theory} --- encoded in the Möbius function.
  \item \textbf{Analytic number theory} --- via the Dirichlet series $1/\zeta(s)$
        and the explicit formula.
  \item \textbf{Probabilistic heuristics} --- the ``random'' cancellation of
        $\mu(n)$ values.
\end{enumerate}

Mertens' three theorems provide the foundational estimates for the distribution
of primes.  The equivalence $\text{RH} \Leftrightarrow M(n)=\BigO(n^{1/2+\varepsilon})$
shows that the Mertens function is not merely a curiosity but a \emph{window}
into the deepest structure of $\zeta$.  The surprise of 1985---that the
intuitively appealing conjecture $|M(n)| < \sqrt{n}$ is false, yet so
elusive that no counterexample has been found---illustrates the subtlety of
analytic number theory.

% ============================================================
\begin{thebibliography}{10}

\bibitem{Mertens1874}
F.~Mertens,
\textit{Ein Beitrag zur analytischen Zahlentheorie},
Crelle's Journal für die reine und angewandte Mathematik \textbf{78} (1874),
46--62.

\bibitem{OdlyzkoTeRiele1985}
A.~M.~Odlyzko and H.~J.~J.~te~Riele,
\textit{Disproof of the Mertens Conjecture},
Journal für die reine und angewandte Mathematik \textbf{357} (1985), 138--160.

\bibitem{Titchmarsh1986}
E.~C.~Titchmarsh (revised by D.~R.~Heath-Brown),
\textit{The Theory of the Riemann Zeta-Function},
2nd ed., Oxford University Press, 1986.

\bibitem{Davenport2000}
H.~Davenport,
\textit{Multiplicative Number Theory},
3rd ed., Springer, 2000.

\bibitem{IwaniecKowalski2004}
H.~Iwaniec and E.~Kowalski,
\textit{Analytic Number Theory},
American Mathematical Society, 2004.

\bibitem{Pintz1987}
J.~Pintz,
\textit{An effective disproof of the Mertens conjecture},
Astérisque \textbf{147--148} (1987), 325--333.

\bibitem{KotnikTeRiele2006}
T.~Kotnik and H.~J.~J.~te~Riele,
\textit{The Mertens Conjecture revisited},
Algorithmic Number Theory (ANTS-VII), Lecture Notes in Computer Science
\textbf{4076} (2006), 156--167.

\bibitem{Selberg1949}
A.~Selberg,
\textit{An Elementary Proof of the Prime Number Theorem},
Annals of Mathematics \textbf{50} (1949), 305--313.

\bibitem{Apostol1976}
T.~M.~Apostol,
\textit{Introduction to Analytic Number Theory},
Springer, 1976.

\bibitem{Haselgrove1958}
C.~B.~Haselgrove,
\textit{A Disproof of a Conjecture of Pólya},
Mathematika \textbf{5} (1958), 141--145.

\end{thebibliography}

\end{document}
